\chapter{Validation, Verification \& Change Management}
\section{Requirements Traceability Matrix}
\par\needspace{8\baselineskip}\begingroup\scriptsize\setlength{\tabcolsep}{4pt}
\begin{longtable}{>{\raggedright\arraybackslash}p{\dimexpr 0.60000\linewidth-2\tabcolsep\relax}>{\raggedright\arraybackslash}p{\dimexpr 0.20000\linewidth-2\tabcolsep\relax}>{\raggedright\arraybackslash}p{\dimexpr 0.20000\linewidth-2\tabcolsep\relax}}
\caption{Functional requirement link coverage}\label{tab:coverage}\\
\toprule
\textbf{Criterion} & \textbf{Count} & \textbf{Coverage} \\ \midrule
\endfirsthead
\multicolumn{3}{l}{\small\itshape Functional requirement link coverage (continued)}\\
\toprule
\textbf{Criterion} & \textbf{Count} & \textbf{Coverage} \\ \midrule
\endhead
\midrule
\multicolumn{3}{r}{\scriptsize Continued on the next page}\\
\endfoot
\bottomrule\endlastfoot
Has a use-case link & 31 & 100\% \\
Has a test-case link & 31 & 100\% \\
Has a direct elicitation-session link & 21 & 68\% \\
\end{longtable}\endgroup

\par\begingroup\scriptsize\setlength{\tabcolsep}{4pt}
\begin{longtable}{>{\raggedright\arraybackslash}p{\dimexpr 0.13000\linewidth-2\tabcolsep\relax}>{\raggedright\arraybackslash}p{\dimexpr 0.27000\linewidth-2\tabcolsep\relax}>{\raggedright\arraybackslash}p{\dimexpr 0.34000\linewidth-2\tabcolsep\relax}>{\raggedright\arraybackslash}p{\dimexpr 0.26000\linewidth-2\tabcolsep\relax}}
\caption{Functional and non-functional traceability}\label{tab:trace}\\
\toprule
\textbf{Requirement} & \textbf{Use case} & \textbf{Test case} & \textbf{HZ / CR / EL} \\ \midrule
\endfirsthead
\multicolumn{4}{l}{\small\itshape Functional and non-functional traceability (continued)}\\
\toprule
\textbf{Requirement} & \textbf{Use case} & \textbf{Test case} & \textbf{HZ / CR / EL} \\ \midrule
\endhead
\midrule
\multicolumn{4}{r}{\scriptsize Continued on the next page}\\
\endfoot
\bottomrule\endlastfoot
REQ-01 & UC-01, UC-07 & TC-01, TC-02, TC-04 & HZ-03; CR-02; EL-01 \\
REQ-02 & UC-01, UC-15 & TC-41, TC-46 & CR-02; EL-01 \\
REQ-03 & UC-02, UC-17 & TC-03, TC-07, TC-26 & HZ-02; CR-01; EL-01 \\
REQ-04 & UC-02 & TC-25 & EL-01 \\
REQ-05 & UC-03, UC-13 & TC-05, TC-27, TC-28 & CR-02; EL-03 \\
REQ-06 & UC-04 & TC-06, TC-29 & EL-03 \\
REQ-07 & UC-06 & TC-10, TC-31, TC-32 & EL-01 \\
REQ-08 & UC-05 & TC-09, TC-15, TC-30, TC-46 & EL-02 \\
REQ-09 & UC-05, UC-10, UC-12, UC-14, UC-16, UC-17, UC-19 & TC-08, TC-20, TC-40, TC-51, TC-55, TC-56 & HZ-05; EL-02 \\
REQ-10 & UC-16 & TC-43 & EL-02 \\
REQ-11 & UC-07, UC-21 & TC-11, TC-12, TC-13 & HZ-06; EL-04 \\
REQ-12 & UC-08, UC-18 & TC-14, TC-15 & EL-05 \\
REQ-13 & UC-09 & TC-16, TC-17 & HZ-08; EL-04 \\
REQ-14 & UC-09 & TC-18 & EL-04 \\
REQ-15 & UC-10, UC-22 & TC-19, TC-20, TC-21, TC-22, TC-44, TC-45, TC-66 & HZ-07; EL-04 \\
REQ-16 & UC-11 & TC-07, TC-23 & HZ-09; CR-01; EL-05 \\
REQ-17 & UC-11 & TC-24 & HZ-09; CR-01; EL-05 \\
REQ-18 & UC-05, UC-12 & TC-33, TC-34, TC-35 & HZ-05; EL-02 \\
REQ-19 & UC-13 & TC-36, TC-37 & EL-03 \\
REQ-20 & UC-14 & TC-38, TC-39, TC-40 & EL-05 \\
REQ-21 & UC-16 & TC-42, TC-47 & EL-05 \\
REQ-22 & UC-17 & TC-50, TC-51 & HZ-02, HZ-12 \\
REQ-23 & UC-18, UC-19 & TC-52, TC-53, TC-54, TC-55, TC-56 & --- \\
REQ-24 & UC-20 & TC-57 & CR-02 \\
REQ-25 & UC-20 & TC-58, TC-59 & CR-02 \\
REQ-26 & UC-02, UC-21 & TC-60, TC-61, TC-62 & HZ-10 \\
REQ-27 & UC-11 & TC-48 & --- \\
REQ-28 & UC-11 & TC-49 & HZ-09 \\
REQ-29 & UC-19, UC-22 & TC-56, TC-63, TC-64, TC-66 & HZ-11 \\
REQ-30 & UC-22 & TC-65 & HZ-11 \\
REQ-31 & UC-01, UC-02 & TC-67 & HZ-04 \\
NFR-01 & --- & TC-04 & EL-02 \\
NFR-02 & --- & TC-67 & HZ-04, HZ-11; EL-02 \\
NFR-03 & --- & --- & --- \\
NFR-04 & --- & --- & HZ-01 \\
NFR-05 & --- & --- & --- \\
NFR-06 & --- & TC-22 & HZ-07; EL-04 \\
\end{longtable}\endgroup

Here, 100\% coverage means that each functional requirement has at least one use-case link and one test-case link. It does not mean every clause is tested, every branch is executed, or stakeholders have approved the requirement. REQ-07, in particular, still lacks acceptance criteria despite having simulation tests.

Dependencies form a directed graph, and a requirement may have multiple parents. The tool displays that graph as a tree, so an identifier can appear more than once. Impact analysis uses sets to avoid duplicate counting. The current dataset version has no general-purpose dependency-cycle validation.

\section{Automated Test Results}
\subsection{Scope and Method}
The runner maps TC-01 through TC-67 to functions in \texttt{checks.py}. Each function operates on separate simulation objects and checks outcomes with assertions. The report calls \texttt{checks.run\_tests(data)} directly rather than invoking workbook write-back. The original data is preserved, and the resulting execution snapshot is stored under \texttt{evidence/}.
\par\needspace{6\baselineskip}\begingroup\footnotesize\setlength{\tabcolsep}{4pt}
\begin{longtable}{>{\raggedright\arraybackslash}p{\dimexpr 0.25000\linewidth-2\tabcolsep\relax}>{\raggedright\arraybackslash}p{\dimexpr 0.25000\linewidth-2\tabcolsep\relax}>{\raggedright\arraybackslash}p{\dimexpr 0.25000\linewidth-2\tabcolsep\relax}>{\raggedright\arraybackslash}p{\dimexpr 0.25000\linewidth-2\tabcolsep\relax}}
\caption{Recorded automated test summary}\\
\toprule
\textbf{Total} & \textbf{Passed} & \textbf{Failed} & \textbf{Not Run} \\ \midrule
\endfirsthead
\multicolumn{4}{l}{\small\itshape Recorded automated test summary (continued)}\\
\toprule
\textbf{Total} & \textbf{Passed} & \textbf{Failed} & \textbf{Not Run} \\ \midrule
\endhead
\midrule
\multicolumn{4}{r}{\scriptsize Continued on the next page}\\
\endfoot
\bottomrule\endlastfoot
67 & 67 & 0 & 0 \\
\end{longtable}\endgroup

The tests cover normal operation, capacity shortages, authorization, fee boundaries, camera failures, sensors, evacuation, monthly/VIP passes, EVs, exit-plate verification, and offline behavior. Important boundaries include the 15/16-minute grace period, 24/25-hour cap, 79/80\% recognition confidence, and the failed-login count that triggers lockout.
Table~\ref{tab:tests} lists every test with its recorded result. NFR-01 and NFR-06 timings quoted there come from the test run and differ slightly from Table~\ref{tab:nfr-results}, which comes from a separate measurement run; both vary with the machine.
\par\begingroup\scriptsize\setlength{\tabcolsep}{4pt}
\begin{longtable}{>{\raggedright\arraybackslash}p{\dimexpr 0.23000\linewidth-2\tabcolsep\relax}>{\raggedright\arraybackslash}p{\dimexpr 0.28000\linewidth-2\tabcolsep\relax}>{\raggedright\arraybackslash}p{\dimexpr 0.49000\linewidth-2\tabcolsep\relax}}
\caption{Complete results of 67 automated tests}\label{tab:tests}\\
\toprule
\textbf{ID / requirements} & \textbf{Test objective} & \textbf{Actual result} \\ \midrule
\endfirsthead
\multicolumn{3}{l}{\small\itshape Complete results of 67 automated tests (continued)}\\
\toprule
\textbf{ID / requirements} & \textbf{Test objective} & \textbf{Actual result} \\ \midrule
\endhead
\midrule
\multicolumn{3}{r}{\scriptsize Continued on the next page}\\
\endfoot
\bottomrule\endlastfoot
TC-01 / REQ-01 & Entry assigns the nearest free slot of the matching type & Passed: Ticket issued; slot A1 assigned. \\
TC-02 / REQ-01 & Entry is rejected when the zone is full & Passed: Entry rejected ('zone full'); rejection logged. \\
TC-03 / REQ-03 & Fee is calculated correctly for a 2h30 stay & Passed: Fee = 30000 VND. \\
TC-04 / REQ-01; NFR-01 & Slot assignment responds within 1 second under normal load & Passed: p95 = 2.332 ms over 1000 requests. \\
TC-05 / REQ-05 & Reserved slot is not given to another vehicle & Passed: Other car got A2; VIP-1 got A1. \\
TC-06 / REQ-06 & Vehicle is located by plate number & Passed: Found: Smart Parking \textgreater{} F1 \textgreater{} A \textgreater{} A2. \\
TC-07 / REQ-03; REQ-16 & Peak-hour rate is applied once CR-01 is implemented & Passed: CR-01 implemented: 17:00-18:00 charged 15000 VND (1.5 x 10000). \\
TC-08 / REQ-09 & Attendant account cannot lock a zone & Passed: PermissionError raised; denial logged. \\
TC-09 / REQ-08 & Locked zone is not assigned to new vehicles & Passed: Car assigned C1 (zone A blocked). \\
TC-10 / REQ-07 & Lost ticket is charged the maximum daily fee plus penalty & Passed: Charged 290000 VND (24h x 10000 + 50000). \\
TC-11 / REQ-11 & Plate with high confidence is accepted automatically & Passed: Plate 29A-111 accepted at 95\% confidence. \\
TC-12 / REQ-11 & Plate with low confidence requires manual entry & Passed: 55\% confidence: manual entry requested and accepted. \\
TC-13 / REQ-11 & Camera offline falls back to manual entry & Passed: Manual entry required; camera\_offline logged. \\
TC-14 / REQ-12 & Nearest free slot is assigned & Passed: Assigned B1 (10 m), not A1 (40 m). \\
TC-15 / REQ-12; REQ-08 & Next nearest slot is assigned when the nearest zone is locked & Passed: Zone B locked: assigned A1. \\
TC-16 / REQ-13 & Sensor mismatch raises alerts & Passed: Alerts: missing\_vehicle, unauthorized\_vehicle. \\
TC-17 / REQ-13 & Silent sensor blocks its slot & Passed: B2 blocked on silence, freed on recovery. \\
TC-18 / REQ-14 & Overstay alert after 24 hours & Passed: 29A-111 parked 25 h in A1 (limit 24 h) \\
TC-19 / REQ-15 & Evacuation opens barriers and stops new entries & Passed: Barriers open; entry rejected; exit charged 0 VND. \\
TC-20 / REQ-15; REQ-09 & Attendant cannot trigger evacuation & Passed: PermissionError raised; denial logged. \\
TC-21 / REQ-15 & Fire alarm starts evacuation without login & Passed: Evacuation started by fire alarm; barriers open. \\
TC-22 / REQ-15; NFR-06 & Barriers open within 5 seconds of the evacuation trigger & Passed: Barriers open in 0.404 ms (500 cars parked). \\
TC-23 / REQ-16 & Peak-hour multiplier applies only to peak hours & Passed: Fee = 25000 VND (10000 + 15000). \\
TC-24 / REQ-17 & Occupancy surcharge applies above 80\% occupancy & Passed: 3rd car base fee; 4th car (100\% occupancy) 12500 VND = 125\% of 10000. \\
TC-25 / REQ-04 & Barrier opens only after payment is made & Passed: Failed payment: slot stays occupied; after payment the slot is free. \\
TC-26 / REQ-03 & Exit with an unknown ticket is refused & Passed: Exit refused: ticket not found. \\
TC-27 / REQ-05 & Reservation is rejected when no slot is free & Passed: 5th reservation rejected (4 car slots reserved). \\
TC-28 / REQ-05 & No-show reservation is released after 15 minutes & Passed: A1 released at 10:46. \\
TC-29 / REQ-06 & Locate rejects an unknown plate & Passed: Plate NOPE: vehicle not found. \\
TC-30 / REQ-08 & Occupied slot stays occupied when its zone is locked & Passed: A1 occupied after lock; blocked after exit. \\
TC-31 / REQ-07 & Lost ticket without an entry record is escalated & Passed: no entry record: escalate to the Parking Manager. \\
TC-32 / REQ-07 & Lost ticket without proof of ownership is refused & Passed: Exit refused; refusal logged. \\
TC-33 / REQ-18 & Valid login returns the account role & Passed: admin1 logged in with role admin. \\
TC-34 / REQ-18 & Wrong password is rejected & Passed: Login rejected; failed attempts: 1. \\
TC-35 / REQ-18 & Account is locked after 3 failed attempts & Passed: account locked (correct password refused after 3 failures). \\
TC-36 / REQ-19 & Cancelling a reservation frees the slot & Passed: A1 released and assigned to the next car. \\
TC-37 / REQ-19 & Cancelling a missing reservation is reported & Passed: Cancel NOPE: reservation not found. \\
TC-38 / REQ-20 & Rate change applies to new tickets only & Passed: Car A pays 10000; car B pays 12000. \\
TC-39 / REQ-20 & Zone added at runtime is assignable at once & Passed: Zone D added; the next car got D1. \\
TC-40 / REQ-20; REQ-09 & Attendant cannot change rates or add zones & Passed: Both attempts denied. \\
TC-41 / REQ-02 & Availability counts match the slot states & Passed: A 0/2, F1 2/4, lot 4/6 free. \\
TC-42 / REQ-21 & Revenue per zone and floor sums to the total & Passed: Zones \{'A': 20000, 'C': 10000\} = floors \{'F1': 20000, 'F2': 10000\} = total 30000. \\
TC-43 / REQ-10 & Audit log records entry, payment and exit & Passed: Audit events: entry, payment, exit. \\
TC-44 / REQ-15 & Ending an evacuation restores normal operation & Passed: Barriers closed, slots unblocked, new entry accepted. \\
TC-45 / REQ-15 & Evacuation cancels reservations & Passed: Reservation cancelled; A1 blocked. \\
TC-46 / REQ-02; REQ-08 & Locked zone shows zero free slots & Passed: Zone A shows 0 free slots. \\
TC-47 / REQ-21 & Report for a period without payments shows zero & Passed: Revenue 0, 0 payments. \\
TC-48 / REQ-27 & A stay of 15 minutes is free & Passed: 15 min = 0 VND; 16 min = 10000 VND. \\
TC-49 / REQ-28 & The fee of a started 24 hours is capped & Passed: 24 h = 100000 VND (cap); 25 h = 110000 VND. \\
TC-50 / REQ-22 & An approved fee override replaces the calculated fee & Passed: ticket fee 5000 VND (was 20000); exit charged 5000; logged: override\_requested, override\_approved \\
TC-51 / REQ-22; REQ-09 & Invalid, unapproved or rejected overrides do not change the fee & Passed: rejected: override needs a non-negative amount and a reason; manager rejects OV-1: fee stays 20000 VND; attendant approval denied: manager or admin only \\
TC-52 / REQ-23 & A monthly pass parks free of charge & Passed: Monthly pass: 5 h stay costs 0 VND. \\
TC-53 / REQ-23 & Only VIP pass holders get VIP slots & Passed: VIP holder: V1; monthly holder: A1; ordinary car: A2. \\
TC-54 / REQ-23 & An expired pass pays the normal fee & Passed: Expired pass: expiry logged, fee 20000 VND. \\
TC-55 / REQ-23; REQ-09 & Administrator registers and removes a pass & Passed: pass stored; logged: pass\_added (monthly until 2026-11-16); MON-1 removed; removal logged \\
TC-56 / REQ-23; REQ-09; REQ-29 & Invalid, unauthorized or offline pass registration is refused & Passed: rejected: invalid pass (kind monthly|vip, end date in the future); invalid pass (kind monthly|vip, end date in the future); attendant denied: administrator only; refused: pass registration needs the central server (offline mode) \\
TC-57 / REQ-24 & EVs get EV slots and cars never do & Passed: EVs got E1 and E2; third EV rejected; the car got A1. \\
TC-58 / REQ-25 & The energy price is added to the fee & Passed: Fee 59000 VND (24000 parking + 35000 energy). \\
TC-59 / REQ-25 & A bad energy reading is rejected & Passed: 0 kWh and a car reading were rejected. \\
TC-60 / REQ-26 & A matching exit plate lets the vehicle leave & Passed: Matching plate: vehicle left. \\
TC-61 / REQ-26 & A plate mismatch keeps the barrier closed & Passed: Mismatch: exit refused, slot still occupied, mismatch logged. \\
TC-62 / REQ-26 & An attendant can approve a mismatch or an unreadable plate & Passed: attendant approves: vehicle leaves (10000 VND), approval logged; unreadable plate: barrier stays closed; a guest cannot approve \\
TC-63 / REQ-29 & Gates keep working offline & Passed: Offline: entry and cash exit worked; 3 records marked. \\
TC-64 / REQ-29 & Online payment and reservations are refused offline & Passed: online refused: payment gateway unavailable: pay cash or card at the booth; slot Occupied; refused offline: reservation needs the central server (offline mode); pass registration needs the central server (offline mode) \\
TC-65 / REQ-30 & Records are synchronized when the network returns & Passed: 1 record synchronized; marks cleared; count logged. \\
TC-66 / REQ-29; REQ-15 & Evacuation works offline & Passed: evacuation offline: barriers OPEN, exit free of charge \\
TC-67 / REQ-31; NFR-02 & Barrier stays open while a vehicle is present or the sensor fails & Passed: Barrier stayed open with a vehicle present, a sensor fault, and power loss (3 refusals logged); closed when clear. \\
\end{longtable}\endgroup


\subsection{Requirements Data Quality Assessment}
\par\needspace{12\baselineskip}\begingroup\footnotesize\setlength{\tabcolsep}{4pt}
\begin{longtable}{>{\raggedright\arraybackslash}p{\dimexpr 0.12000\linewidth-2\tabcolsep\relax}>{\raggedright\arraybackslash}p{\dimexpr 0.16000\linewidth-2\tabcolsep\relax}>{\raggedright\arraybackslash}p{\dimexpr 0.72000\linewidth-2\tabcolsep\relax}}
\caption{Requirements data quality findings}\label{tab:quality}\\
\toprule
\textbf{Level} & \textbf{ID} & \textbf{Finding} \\ \midrule
\endfirsthead
\multicolumn{3}{l}{\small\itshape Requirements data quality findings (continued)}\\
\toprule
\textbf{Level} & \textbf{ID} & \textbf{Finding} \\ \midrule
\endhead
\midrule
\multicolumn{3}{r}{\scriptsize Continued on the next page}\\
\endfoot
\bottomrule\endlastfoot
ERROR & REQ-07 & missing acceptance criteria \\
ERROR & NFR-05 & not measurable: metric or target value missing \\
WARN & REQ-07 & not atomic: description joins several actions (split into separate requirements) \\
WARN & REQ-10 & not atomic: description joins several actions (split into separate requirements) \\
INFO & REQ-07 & Must requirement still Proposed (needs approval) \\
INFO & REQ-10 & Must requirement still Proposed (needs approval) \\
INFO & REQ-15 & Must requirement still Proposed (needs approval) \\
\end{longtable}\endgroup

The two ERROR findings concern specification and measurement data; they remain even when every test is Passed. Atomicity warnings use a comma-and-conjunction heuristic and require analyst review. INFO findings identify Must requirements that remain Proposed, so business approval is needed before establishing an official baseline.

Further specification/implementation gaps include the lost-ticket rule versus the daily cap, calls to \texttt{exit} without an observed plate, zone locks before and after evacuation, and persistence of dynamically added zones. Reporting these gaps does not change the project's source code.

\section{Requirements Change Management}
The proposed process records a CR, identifies directly affected requirements, calculates dependency closure, reviews use cases/tests/risks, estimates effort, obtains a decision, updates the baseline, and runs regression tests. The status \emph{Implemented} in the workbook means the change was applied in the simulation, not accepted on a real installation.
\par\needspace{7\baselineskip}\begingroup\footnotesize\setlength{\tabcolsep}{4pt}
\begin{longtable}{>{\raggedright\arraybackslash}p{\dimexpr 0.12000\linewidth-2\tabcolsep\relax}>{\raggedright\arraybackslash}p{\dimexpr 0.20000\linewidth-2\tabcolsep\relax}>{\raggedright\arraybackslash}p{\dimexpr 0.52000\linewidth-2\tabcolsep\relax}>{\raggedright\arraybackslash}p{\dimexpr 0.16000\linewidth-2\tabcolsep\relax}}
\caption{Impact summary for the two change requests}\label{tab:changes}\\
\toprule
\textbf{CR} & \textbf{Direct requirements} & \textbf{Impact closure (requirements, use cases, tests)} & \textbf{Sample estimate} \\ \midrule
\endfirsthead
\multicolumn{4}{l}{\small\itshape Impact summary for the two change requests (continued)}\\
\toprule
\textbf{CR} & \textbf{Direct requirements} & \textbf{Impact closure (requirements, use cases, tests)} & \textbf{Sample estimate} \\ \midrule
\endhead
\midrule
\multicolumn{4}{r}{\scriptsize Continued on the next page}\\
\endfoot
\bottomrule\endlastfoot
CR-01 / Implemented & REQ-03, REQ-16, REQ-17 & 16 requirements (REQ-03, REQ-04, REQ-07, REQ-10, REQ-16, REQ-17, REQ-20, REQ-21, REQ-22, REQ-23, REQ-25, REQ-26, REQ-27, REQ-28, REQ-29, REQ-30); 11 use cases; 33 tests & cost 5 days; schedule 7 days \\
CR-02 / Implemented & REQ-01, REQ-02, REQ-05, REQ-24, REQ-25 & 29 requirements (REQ-01, REQ-02, REQ-03, REQ-04, REQ-05, REQ-06, REQ-07, REQ-08, REQ-10, REQ-11, REQ-12, REQ-13, REQ-14, REQ-15, REQ-16, REQ-17, REQ-19, REQ-20, REQ-21, REQ-22, REQ-23, REQ-24, REQ-25, REQ-26, REQ-27, REQ-28, REQ-29, REQ-30, REQ-31); 21 use cases; 63 tests & cost 3 days; schedule 3 days \\
\end{longtable}\endgroup

CR-01 adds peak-hour and occupancy pricing. Regression should include grace periods, passes, and revenue recording because they share fee calculation. CR-02 adds EV allocation and charging and has broad impact through the original allocation dependency. The returned impact set is a conservative review aid; final regression priorities and scope still require business interpretation.

\section{Chapter Conclusion}
The recorded execution shows the existing suite passing on the simulation, while traceability identifies the scope of changes. Specification quality, approval, and behavioral coverage require separate assessment. A passing-test percentage cannot replace requirements validation.

The suite contains 67 automated tests and all 67 pass; every functional requirement links to at least one use case and one test (Table~\ref{tab:coverage}). The quality checker still reports two errors, two warnings, and three informational findings, which no test can remove, because they concern the wording of requirements and not the behavior of the code. The impact of the two change requests was computed through the dependency tree (Table~\ref{tab:changes}), and both are marked Implemented in the simulation. That status does not mean acceptance on a real installation. Validation by real stakeholders is the missing step.
